\documentclass[12pt, french, latin.medieval, a4paper]{report}
% package de langues 
\usepackage[french]{babel}
\usepackage[utf8]{inputenc}
\usepackage[T1]{fontenc}

% package pour référence de liens sur pdf 
\usepackage[pdftex]{hyperref} % liens
\usepackage{lastpage} % numéro de page
\usepackage{caption}

% package de formes de document 
\usepackage[left=2.5cm,
              right=2.5cm,
              top=2.5cm,
              bottom=2.5cm
             ]{geometry}
% pour les tableaux 
\usepackage{array,multirow,makecell} 
\newcolumntype{R}[1]{>{\raggedleft\arraybackslash }b{#1}}
\newcolumntype{L}[1]{>{\raggedright\arraybackslash }b{#1}}
\newcolumntype{C}[1]{>{\centering\arraybackslash }b{#1}}

% package pour la page de couverture 
\usepackage{graphicx}
\usepackage[cyr]{aeguill}
\usepackage{fancyheadings}
\usepackage{amsmath}
\usepackage{amsfonts}
\usepackage{amssymb}
\usepackage{geometry}
\usepackage{eso-pic}
\usepackage{tikz}
\usetikzlibrary{calc}


% nouvelle commande pour la forme 
%\fancyhf[OFC]{\thepage / \pageref{LastPage}}
\rfoot{\thepage / \pageref{LastPage}}
%\setcounter{tocdepth}{3}

% pakages de couleurs du document
\usepackage{color}
\usepackage{xcolor}


% package pour la chimie
\usepackage[version=4]{mhchem}


% package pour la visualisation du code
\usepackage{listings}
\definecolor{darkWhite}{rgb}{0.94,0.94,0.94}
 
\lstset{
  aboveskip=3mm,
  belowskip=-2mm,
  backgroundcolor=\color{darkWhite},
  basicstyle=\footnotesize,
  breakatwhitespace=false,
  breaklines=true,
  captionpos=b,
  commentstyle=\color{red},
  deletekeywords={...},
  escapeinside={\%*}{*)},
  extendedchars=true,
  framexleftmargin=16pt,
  framextopmargin=3pt,
  framexbottommargin=6pt,
  frame=tb,
  keepspaces=true,
  keywordstyle=\color{blue},
  language=python,
  literate=
  {²}{{\textsuperscript{2}}}1
  {⁴}{{\textsuperscript{4}}}1
  {⁶}{{\textsuperscript{6}}}1
  {⁸}{{\textsuperscript{8}}}1
  {€}{{\euro{}}}1
  {é}{{\'e}}1
  {è}{{\`{e}}}1
  {ê}{{\^{e}}}1
  {ë}{{\¨{e}}}1
  {É}{{\'{E}}}1
  {Ê}{{\^{E}}}1
  {û}{{\^{u}}}1
  {ù}{{\`{u}}}1
  {â}{{\^{a}}}1
  {à}{{\`{a}}}1
  {á}{{\'{a}}}1
  {ã}{{\~{a}}}1
  {Á}{{\'{A}}}1
  {Â}{{\^{A}}}1
  {Ã}{{\~{A}}}1
  {ç}{{\c{c}}}1
  {Ç}{{\c{C}}}1
  {õ}{{\~{o}}}1
  {ó}{{\'{o}}}1
  {ô}{{\^{o}}}1
  {Õ}{{\~{O}}}1
  {Ó}{{\'{O}}}1
  {Ô}{{\^{O}}}1
  {î}{{\^{i}}}1
  {Î}{{\^{I}}}1
  {í}{{\'{i}}}1
  {Í}{{\~{Í}}}1,
  morekeywords={*,...},
  numbers=left,
  numbersep=10pt,
  numberstyle=\tiny\color{black},
  rulecolor=\color{black},
  showspaces=false,
  showstringspaces=false,
  showtabs=false,
  stepnumber=1,
  stringstyle=\color{gray},
  tabsize=4,
  title=\lstname,
}

%package pour les algorithmes
\usepackage{algorithm}
\usepackage{algorithmic}
\usepackage{amsmath}

% package pour les tableaux
\usepackage{longtable}
\usepackage{booktabs}

% package pour les figures 
\usepackage{subcaption}

% package de reférencement
% \usepackage[backend=biber,safeinputenc, sorting=none]{biblatex}

% mettre la bibliographie dans la table de matières
%\usepackage[nottoc, notlof, notlot]{tocbibind}
\include{page_couverture.tex}


\author{MBE MBE MINDJANA LOIC HENRI}
\title{GNPS et MDByTE}

\begin{document}
    
    \pagecouverture
    {Idée de vote majoritaire et pond ́er ́e pour les mod`eles d’apprentissage
par ensemble d’algorithmes de clustering}
    {}
	\tableofcontents
	\vspace{20cm}
    
    \section{Introduction}
    Dans le but de combiner GNPS et MADByTE à des algorithmes de clustering multi-vues, pour tenir compte de la complémentarité des spectres MS et RMN, on pourrait définir un modèle d’apprentissage par ensemble.
Ce modèle doit pouvoir être capable de combiner des algorithmes de clustering fournissant un nombre fini de clusters ou non, pour obtenir un nombre fini de clusters.
    	Pour le faire, il est important de définir comment le vote pondéré et le vote majoritaire vont se réaliser.
   	A travers ce document, nous proposons un moyen de vote pondéré pour un modèle d'apprentissage par ensemble d'algorithmes de clustering de façon générale.

    \section{Vote pondérée et majoritaire}
   	 Considérons que pour caractériser les clusters finaux, chaque observation $i$ possède des poids pour chaque cluster $k$. 
Notons $p_{i,:}$ le vecteur de ses poids qu’on va considérer comme son indicateur de cluster, où la $k$-eme composante correspond au poids pour le $k$-eme cluster.
   	 
   	 On peut considérer le fait qu'un weak-learner $w$ valide l'attribution d’indicateurs à deux observations $i$ et $j$ qui sont dans un même groupe (pour lui), si leurs indicateurs sont semblables.
   	 De ce fait, il sera moins satisfait si la différence de ces indicateurs est grande, surtout s'il a un grande importance de vote. Ainsi le vote pour obtenir les indicateurs de deux observations peut se faire, en réduisant le plus possible la non satisfaction des weak-learners qui ont de grandes importances comparé à ceux qui en ont moins.

   	 De là on peut quantifier la non satisfaction générale d'attribuer des indicateurs à deux observations $i$ et $j$, par $nS_{i,j}(P)$ défini comme suit, où $\lambda(w)$ correspond au poids d’importance du weak-learner $w$ pour le vote :
   	 \begin{equation}
   		 \begin{aligned}
   			\\ & nS_{i,j}(P) = \sum_{w} \lambda(w) {\Vert p_{i,:} - p_{j,:} \Vert}_2^2 A_{i,j}(w)
            \\ & A_{i,j}(w) = \begin{cases}
                    1 & \text{ si $i \neq j$, l'observation $i$ et $j$ sont du même groupe d'apres le weak-learner $w$}  \\
                    0 & \text{sinon} \\
                \end{cases}
   		 \end{aligned}
   	 \end{equation}
   	 D’où on a la non satisfaction globale $NS(P)$, étant:
   	 \begin{equation}
   		 \begin{aligned}
   			 & NS(P) = \sum_{i,j,w} \lambda(w) {\Vert p_{i,:} - p_{j,:} \Vert}_2^2 A_{i,j}(w)
   		 \end{aligned}
   	 \end{equation}

   	 Sachant qu'on souhaite avoir la plus petite possible non satisfaction globale $NS(P)$, et que les clusters ne doivent pas avoir de corrélation entre elle, on a le problème d'optimisation suivant à résoudre, où $n$ est le nombre d’observations et $k$ le nombre de clusters voulu :
   	 \begin{equation}
   		 \begin{aligned}
   			 & \min_{P} \sum_{i,j,w} \lambda(w) {\Vert p_{i,:} - p_{j,:} \Vert}_2^2 A_{i,j}(w)
   			 \\ & P \in \mathbb{R}^{n \times k} ,\ P^TP = I
   		 \end{aligned}
   	 \end{equation}
     
   	 Notons que la contrainte $P^TP = I$ permet de vérifier la non colinéarité entre deux différents clusters $k1$ (de vecteurs de poids $P_{:,k1}$) et $k2$ (de vecteurs de poids $P_{:,k2}$), car implique que $<P_{:,k1}, P_{:,k2}>=0$,
   	 et permet également d'indiquer en bonus que les vecteurs de poids des clusters seront unitaires.
   	 
   	 En utilisant les conditions de KKT, on peut montrer que résoudre le problème revient à trouver les $k$ premiers vecteurs propres (correspondant aux vecteurs de poids des clusters) de la matrice $L$ défini comme suit:
     \begin{equation}
   		 \begin{aligned}
   			L = \sum_w D(w) - A(w)
            \\ D_{i,j}(w) = \begin{cases}
                    \sum_j A_{i,j}(w) & \text{ si $i = j$}  \\
                    0 & \text{sinon} \\
                \end{cases}
   		 \end{aligned}
   	 \end{equation}

   	 Après obtention de $P$, on pourra obtenir les clusters proprement dit, en utilisant l'algorithme KMEANS sur $P$, pour plusieurs  raisons:
   	 \begin{itemize}
   		 \item P correspond à une sorte de projection des observations l’espace des clusters, dont il est fort probable qu’on retrouve des groupes sphériques.
   		 \item pour cette raison, il est le plus adéquat, car il est simple, rapide et efficace des algorithmes de clustering sur des groupes sphériques à nombre voulu.
   		 \item il a été recommandé et expérimenté dans les articles de référence d'algorithme tel que MCLES, des algorithmes de clustering spectral multivues et à vue unique, et utiliser dans leur implémentation respectif
   	 \end{itemize}

   	 De façon analogue, on a le vote majoritaire, en considérant que les weak-learners ont des importances identiques pour le vote.

    \section{Conclusion}
    En somme, il était question pour nous de proposer un moyen de réaliser le vote pondéré et le vote majoritaire des clusters à nombre variable des weak-learners.
    Nous pensons que si les modèles sont performants, on aura pas forcément de gain de performance.
    Il serait important de définir des bons poids d'importance aux weak-learners, car si on donne plus d'importance à ceux moins bons, on aura forcément de mauvaises performances.
	De ce fait, comment définir le fait qu’une observation a été mis dans un bon groupe?
   
\bibliographystyle{plain}
	\bibliography{reference_GNPS_MADByTE}
\end{document}
